\documentclass[11pt]{article}

% ============================================================
% Packages
% ============================================================
\usepackage[margin=1in]{geometry}
\usepackage{amsmath,amssymb,mathtools,amsthm}
\usepackage{enumitem}
\usepackage{hyperref}
\usepackage{xcolor}
\usepackage{microtype}
\usepackage{graphicx}

% ============================================================
% Macros
% ============================================================
\newcommand{\R}{\mathbb{R}}
\newcommand{\Pbb}{\mathbb{P}}
\newcommand{\F}{\mathcal{F}}
\newcommand{\Qbb}{\mathbb{Q}}
\newcommand{\E}{\mathbb{E}}
\newcommand{\dd}{\,\mathrm{d}}

% ============================================================
% Environments
% ============================================================
\theoremstyle{definition}
\newtheorem{definition}{Definition}
\newtheorem{remark}{Remark}

\theoremstyle{plain}
\newtheorem{proposition}{Proposition}
\newtheorem{theorem}{Theorem}

% ============================================================
% Title
% ============================================================
\title{\textbf{Change of Numeraire}\\
\large MSFE6233 --- Option Pricing}
\author{Nathan Sauldubois}
\date{}

\begin{document}

\maketitle

\vspace{-1em}

\section{Motivation and general principle}


In the Black--Scholes framework, and more generally in arbitrage-free
continuous-time models, the standard pricing principle says that
discounted asset prices are martingales under a suitable risk-neutral
measure.
Usually, the discounting is performed using the money market account,
which leads to the familiar statement that
$
D(t)S(t)
$
is a martingale under the risk-neutral measure.

However, there is nothing fundamentally special about the money market
account.
More generally, one may choose any strictly positive traded asset
$
N(t)
$
as a unit of account.
Such an asset is called a \emph{numeraire}.
Once a numeraire is chosen, every asset price can be expressed in units
of this numeraire, i.e.
$
S^N(t):=\frac{S(t)}{N(t)}.
$

The key idea of the change of numeraire is that a change of unit of
account naturally comes with a change of probability measure.
More precisely, for each admissible numeraire \(N\), there exists an
associated equivalent martingale measure, under which prices expressed
in units of \(N\) are martingales.
Thus, instead of working with
$
D(t)S(t),
$
one may work with
$
S(t)/N(t),
$
provided one also changes the pricing measure accordingly.

This point of view is conceptually important and practically very
useful.

\paragraph{Why is this useful?}

First, it gives a more flexible formulation of risk-neutral valuation.
The standard money-market numeraire is only one possible choice among
many others.
Depending on the derivative being priced, another numeraire may lead to
a simpler pricing formula or a more convenient dynamics.

Second, it allows one to simplify the treatment of stochastic interest
rates.
When pricing under the money-market numeraire, one often has to handle
the joint randomness of the payoff and the discount factor.
By choosing a more suitable numeraire, such as a zero-coupon bond
maturing at the payoff date, one can often absorb the discounting into
the numeraire and obtain a cleaner expectation formula.

Third, the change of numeraire is particularly important in fixed-income
and FX modeling.
For instance, choosing a zero-coupon bond as numeraire leads to the
\(T\)-forward measure, under which forward prices become martingales.
Likewise, in FX markets, changing from the domestic to the foreign
money-market account naturally leads to the foreign pricing measure.

\paragraph{General principle.}

The guiding principle of the chapter is therefore the following:

\begin{quote}
Choose a positive traded asset \(N\) as numeraire, express all prices in
units of \(N\), and switch to the associated probability measure under
which the denominated prices become martingales.
\end{quote}

This is exactly the same philosophy as ordinary risk-neutral pricing,
but formulated in a more general and more powerful way.

\paragraph{What will be studied in this lecture?}

We proceed in four steps.

First, we recall the notion of numeraire and explain how a new
probability measure is associated with each choice of numeraire.

Second, we derive the general pricing formula under an arbitrary
numeraire and explain the martingale property of denominated prices.

Third, we revisit It\^o calculus for ratio processes, in order to see
how asset dynamics transform when the numeraire is changed.

Finally, we study the most important examples:
the money market account, the stock as numeraire, and the zero-coupon
bond numeraire, with particular emphasis on forward measures and their
applications to derivative pricing.

\subsection{It\^o calculus refresher for ratio processes}

To use the change-of-numeraire method effectively, it is useful to
recall how It\^o's formula applies to ratios of stochastic processes.

\paragraph{Why ratios appear.}
Once a numeraire \(N(t)\) has been chosen, the relevant price process is
no longer \(S(t)\) itself, but the denominated price
\[
\frac{S(t)}{N(t)}.
\]
The dynamics of this ratio determine whether the process is a martingale
under the associated measure.

\paragraph{General setting.}
Let \(S(t)\) and \(N(t)\) be two strictly positive It\^o processes of the form
\[
\mathrm{d}S(t)=\mu_S(t)S(t)\,\mathrm{d}t+\sigma_S(t)S(t)\,\mathrm{d}W(t),
\]
\[
\mathrm{d}N(t)=\mu_N(t)N(t)\,\mathrm{d}t+\sigma_N(t)N(t)\,\mathrm{d}W(t),
\]
where, for simplicity, we write the one-dimensional case.

We want to compute the dynamics of
\[
R(t):=\frac{S(t)}{N(t)}.
\]

\paragraph{It\^o formula for the ratio.}
Applying It\^o's formula to the map
\[
f(s,n)=\frac{s}{n},
\]
we obtain
\[
\mathrm{d}\Big(\frac{S(t)}{N(t)}\Big)
=
\frac{1}{N(t)}\,\mathrm{d}S(t)
-
\frac{S(t)}{N(t)^2}\,\mathrm{d}N(t)
+
\frac{S(t)}{N(t)^3}\,\mathrm{d}\langle N\rangle_t
-
\frac{1}{N(t)^2}\,\mathrm{d}\langle S,N\rangle_t.
\]

Using
\[
\mathrm{d}\langle N\rangle_t=\sigma_N(t)^2N(t)^2\,\mathrm{d}t,
\qquad
\mathrm{d}\langle S,N\rangle_t=\sigma_S(t)\sigma_N(t)S(t)N(t)\,\mathrm{d}t,
\]
this yields
\[
\mathrm{d}\Big(\frac{S(t)}{N(t)}\Big)
=
\frac{S(t)}{N(t)}
\Big(
\mu_S(t)-\mu_N(t)+\sigma_N(t)^2-\sigma_S(t)\sigma_N(t)
\Big)\mathrm{d}t
+
\frac{S(t)}{N(t)}\big(\sigma_S(t)-\sigma_N(t)\big)\mathrm{d}W(t).
\]

\paragraph{Interpretation of the drift term.}
This formula shows that the drift of the ratio is not simply the
difference of the two drifts.
There is an additional quadratic-variation correction term coming from
It\^o's formula.

This is exactly why a change of numeraire generally requires a change of
measure: under the original measure, the process \(S(t)/N(t)\) does not
usually have zero drift.

\paragraph{Special case: money market account.}
If \(N(t)=B(t)\) is the money market account, then
\[
\mathrm{d}B(t)=r(t)B(t)\,\mathrm{d}t,
\]
so \(\sigma_N(t)=0\).
In that case,
\[
\mathrm{d}\Big(\frac{S(t)}{B(t)}\Big)
=
\frac{S(t)}{B(t)}\big(\mu_S(t)-r(t)\big)\mathrm{d}t
+
\frac{S(t)}{B(t)}\sigma_S(t)\mathrm{d}W(t).
\]

Under the usual risk-neutral measure, we have \(\mu_S(t)=r(t)\), and
therefore
\[
\mathrm{d}\Big(\frac{S(t)}{B(t)}\Big)
=
\frac{S(t)}{B(t)}\sigma_S(t)\mathrm{d}W(t),
\]
which is the familiar martingale property of discounted prices.

\paragraph{Logarithmic viewpoint.}
It is also useful to write
\[
\log\Big(\frac{S(t)}{N(t)}\Big)=\log S(t)-\log N(t).
\]
This can simplify computations when both \(S\) and \(N\) are geometric
It\^o processes.

\paragraph{Main lesson.}
The change of numeraire amounts to studying ratio processes.
It\^o's formula shows explicitly how their drift and volatility are
modified, and it explains why the appropriate martingale property only
holds after the corresponding change of measure.

\paragraph{What comes next.}
We now turn to the main examples of numeraires and illustrate how the
general pricing formula simplifies in concrete situations.


\section{Numeraire, measure, and martingale property}

\subsection{Choice of numeraire}

We begin by formalizing the notion of numeraire.

\paragraph{Definition.}
A \emph{numeraire} is a strictly positive traded asset chosen as the
unit of account.
Once a numeraire \(N(t)\) has been selected, any asset price \(S(t)\)
can be expressed in units of this numeraire through the ratio
\[
S^N(t):=\frac{S(t)}{N(t)}.
\]
The main idea is that pricing can be carried out using these
denominated prices instead of the original dollar prices.  

\paragraph{The standard choice.}
In the usual Black--Scholes framework, the natural numeraire is the
money market account
\[
B(t)=e^{\int_0^t R(s)\,\mathrm{d}s},
\]
or, in the constant-rate case,
\[
B(t)=e^{rt}B(0).
\]
Choosing \(B(t)\) as numeraire leads to the standard discounted price
process
\[
\frac{S(t)}{B(t)}=D(t)S(t),
\qquad D(t)=\frac1{B(t)},
\]
which is a martingale under the usual risk-neutral measure. 

\paragraph{Why can we choose another numeraire?}
The important point is that the money market account is only one
possible choice.
In any arbitrage-free model without dividends, one may use any strictly
positive traded asset as numeraire.
Possible choices include:
\begin{itemize}[leftmargin=1.4em]
\item the domestic money market account,
\item the foreign money market account,
\item a zero-coupon bond maturing at some fixed date \(T\),
\item or even a derivative asset, provided its price remains strictly positive.
\end{itemize}
Each such choice comes with its own associated equivalent martingale
measure. 
\paragraph{Dynamics of the numeraire under the risk-neutral measure.}

Let \(N(t)\) be a strictly positive traded asset, and assume that we work
under the usual risk-neutral measure \(\widetilde{\mathbb P}\)
associated with the money market account
\[
B(t)=\exp\!\Big(\int_0^t R(s)\,\mathrm{d}s\Big).
\]

Since \(N\) is a traded asset, its discounted price
\[
D(t)N(t)=\frac{N(t)}{B(t)}
\]
must be a \(\widetilde{\mathbb P}\)-martingale.

Under the usual square-integrability assumptions, the martingale
representation theorem implies that there exists an adapted process
\(\nu(t)\in\R^d\) such that
\[
\mathrm{d}\big(D(t)N(t)\big)
=
D(t)N(t)\,\nu(t)\cdot \mathrm{d}\widetilde W(t).
\]

Equivalently,
\[
\frac{\mathrm{d}\big(D(t)N(t)\big)}{D(t)N(t)}
=
\nu(t)\cdot \mathrm{d}\widetilde W(t).
\]

Since
\[
\mathrm{d}D(t)=-R(t)D(t)\,\mathrm{d}t,
\]
It\^o's product rule yields
\[
\mathrm{d}\big(D(t)N(t)\big)
=
D(t)\,\mathrm{d}N(t)+N(t)\,\mathrm{d}D(t),
\]
because \(D\) has finite variation and therefore
\[
\mathrm{d}\langle D,N\rangle_t=0.
\]

Substituting the dynamics of \(D\), we obtain
\[
\mathrm{d}\big(D(t)N(t)\big)
=
D(t)\Big(\mathrm{d}N(t)-R(t)N(t)\,\mathrm{d}t\Big).
\]

Comparing this with
\[
\mathrm{d}\big(D(t)N(t)\big)
=
D(t)N(t)\,\nu(t)\cdot \mathrm{d}\widetilde W(t),
\]
we deduce that the numeraire \(N\) satisfies
\[
\boxed{
\mathrm{d}N(t)
=
R(t)N(t)\,\mathrm{d}t
+
N(t)\,\nu(t)\cdot \mathrm{d}\widetilde W(t).
}
\]

\paragraph{Interpretation.}

This formula shows that, under the risk-neutral measure, every strictly
positive traded numeraire has drift equal to the short rate \(R(t)\).
Only its volatility process \(\nu(t)\) depends on the specific choice of
numeraire.

Thus:
\begin{itemize}[leftmargin=1.4em]
\item the money market account corresponds to the special case
\[
\nu(t)=0,
\]
\item a risky numeraire has the same drift \(R(t)\), but a non-zero
volatility process \(\nu(t)\).
\end{itemize}

\paragraph{Integrated form.}

The previous SDE can be written explicitly as
\[
N(t)
=
N(0)\exp\!\Big(
\int_0^t \Big(R(s)-\frac12\|\nu(s)\|^2\Big)\,\mathrm{d}s
+
\int_0^t \nu(s)\cdot \mathrm{d}\widetilde W(s)
\Big).
\]

This makes it clear that \(N(t)\) remains strictly positive, as required
for a valid numeraire.


\paragraph{Main principle.}
Once a numeraire \(N\) is chosen, the relevant object is no longer the
discounted price \(D(t)S(t)\), but rather the denominated price
\[
\frac{S(t)}{N(t)}.
\]
The fundamental theorem is then reformulated as follows: associated with
each numeraire \(N\), there exists an equivalent martingale measure
\(\Qbb^N\) under which prices expressed in units of \(N\) are
martingales. In this sense, the ordinary risk-neutral measure is simply
the special case corresponding to the money market account. 

\paragraph{Dynamics of denominated asset prices under the new measure.}

Let \(S(t)\) be a traded asset and let \(N(t)\) be a strictly positive
traded numeraire.
Under the usual risk-neutral measure \(\widetilde{\mathbb P}\), suppose
that
\[
\mathrm{d}S(t)
=
R(t)S(t)\,\mathrm{d}t
+
S(t)\sigma(t)\cdot \mathrm{d}\widetilde W(t),
\]
and, as shown above,
\[
\mathrm{d}N(t)
=
R(t)N(t)\,\mathrm{d}t
+
N(t)\nu(t)\cdot \mathrm{d}\widetilde W(t).
\]

We define the denominated price process
\[
S^N(t):=\frac{S(t)}{N(t)}.
\]

\paragraph{Step 1: dynamics under \(\widetilde{\mathbb P}\).}

Applying It\^o's formula to the ratio \(S(t)/N(t)\), we obtain
\[
\mathrm{d}S^N(t)
=
S^N(t)
\Big(
\|\nu(t)\|^2-\sigma(t)\cdot\nu(t)
\Big)\mathrm{d}t
+
S^N(t)\big(\sigma(t)-\nu(t)\big)\cdot \mathrm{d}\widetilde W(t).
\]

Thus, under the original risk-neutral measure
\(\widetilde{\mathbb P}\), the denominated price process generally has
a non-zero drift.

\paragraph{Step 2: change of measure.}

Let \(\widetilde{\mathbb P}^N\) be the probability measure associated
with the numeraire \(N\).
Its density with respect to \(\widetilde{\mathbb P}\) is generated by
the process \(\nu\), and Girsanov's theorem implies that
\[
\widetilde W^N(t)
:=
\widetilde W(t)-\int_0^t \nu(s)\,\mathrm{d}s
\]
is a Brownian motion under \(\widetilde{\mathbb P}^N\).

Equivalently,
\[
\mathrm{d}\widetilde W(t)
=
\mathrm{d}\widetilde W^N(t)+\nu(t)\,\mathrm{d}t.
\]

Substituting this into the previous dynamics gives
\begin{align*}
\mathrm{d}S^N(t)
&=
S^N(t)
\Big(
\|\nu(t)\|^2-\sigma(t)\cdot\nu(t)
\Big)\mathrm{d}t
+
S^N(t)\big(\sigma(t)-\nu(t)\big)\cdot
\big(\mathrm{d}\widetilde W^N(t)+\nu(t)\,\mathrm{d}t\big)
\\
&=
S^N(t)\big(\sigma(t)-\nu(t)\big)\cdot \mathrm{d}\widetilde W^N(t),
\end{align*}
since the drift terms cancel.

Hence, under the measure associated with the numeraire \(N\), the
denominated asset price is driftless:
\[
\boxed{
\mathrm{d}S^N(t)
=
S^N(t)\big(\sigma(t)-\nu(t)\big)\cdot \mathrm{d}\widetilde W^N(t).
}
\]

\paragraph{Main consequence.}

This formula makes the change-of-numeraire principle completely explicit:
under the measure associated with \(N\), the denominated price process
\(S(t)/N(t)\) is a local martingale, and in the present setting a true
martingale.

\paragraph{Pedagogical message.}

The important point is the following:
\begin{quote}
changing numeraire removes the drift of the denominated price, but it
also changes its volatility.
\end{quote}

More precisely, the volatility of \(S^N(t)\) under the new measure is
not \(\sigma(t)\), but
\[
\sigma(t)-\nu(t).
\]

Thus a change of numeraire does not simply relabel prices:
it also changes the stochastic dynamics of the relevant processes.

\paragraph{Special case.}

If the numeraire is the money market account, then \(\nu(t)=0\), so
\[
S^N(t)=\frac{S(t)}{B(t)}
\]
and the above formula reduces to the familiar risk-neutral dynamics
\[
\mathrm{d}\Big(\frac{S(t)}{B(t)}\Big)
=
\frac{S(t)}{B(t)}\sigma(t)\cdot \mathrm{d}\widetilde W(t).
\]


\paragraph{Why is the choice of numeraire useful?}
The choice of numeraire is largely a matter of convenience.
A suitable numeraire can simplify both the pricing formula and the asset
dynamics.
For example, choosing a zero-coupon bond maturing at \(T\) leads to the
\(T\)-forward measure, under which the forward price process becomes a
martingale, which is particularly convenient for bond options and
interest-rate derivatives. More generally, different numeraires are
adapted to different classes of problems, such as fixed-income or FX
pricing. 

\paragraph{What we do next.}
In the next subsection, we explain how a probability measure is attached
to a given numeraire and why, under that new measure, the denominated
asset prices become martingales.


\subsection{Associated probability measures}

We now explain how a new probability measure is associated with a given
choice of numeraire.

\paragraph{General idea.}
Suppose that \(N(t)\) is a strictly positive traded asset chosen as
numeraire.
The goal is to construct a probability measure \(\Qbb^N\) under which
all asset prices expressed in units of \(N\) become martingales.

This is the natural extension of the usual risk-neutral construction:
when the money market account \(B(t)\) is chosen as numeraire, the
associated measure is precisely the standard risk-neutral measure.

\paragraph{Radon--Nikodym density.}
Let \(B(t)\) denote the money market account, and let \(\Qbb\) be the
usual risk-neutral measure associated with \(B\).
The probability measure \(\Qbb^N\) associated with the numeraire \(N\)
is defined through the density process
\[
L_t
:=
\frac{\mathrm{d}\Qbb^N}{\mathrm{d}\Qbb}\Big|_{\mathcal{F}_t}
=
\frac{N(t)/B(t)}{N(0)/B(0)}.
\]

Equivalently,
\[
L_t
=
\frac{D(t)N(t)}{D(0)N(0)},
\qquad D(t)=\frac1{B(t)}.
\]

Since \(N\) is a traded asset, the discounted process \(D(t)N(t)\) is a
\(\Qbb\)-martingale.
Therefore \(L_t\) is a positive \(\Qbb\)-martingale with
\[
\E^\Qbb[L_t]=1,
\]
so it indeed defines a new probability measure equivalent to \(\Qbb\).

\paragraph{Change of measure between two numeraires.}

Let \(N(t)\) and \(M(t)\) be two strictly positive traded numeraires, and
let \(\Qbb^N\) and \(\Qbb^M\) be their associated probability measures.

By definition, with respect to the usual risk-neutral measure
\(\widetilde{\mathbb P}\), we have
\[
\frac{\mathrm{d}\Qbb^N}{\mathrm{d}\widetilde{\mathbb P}}\Big|_{\F_t}
=
\frac{D(t)N(t)}{D(0)N(0)},
\qquad
\frac{\mathrm{d}\Qbb^M}{\mathrm{d}\widetilde{\mathbb P}}\Big|_{\F_t}
=
\frac{D(t)M(t)}{D(0)M(0)},
\]
where
\[
D(t)=\frac1{B(t)}.
\]

Taking the ratio of these two Radon--Nikodym derivatives, we obtain the
density process of \(\Qbb^M\) with respect to \(\Qbb^N\):
\[
\frac{\mathrm{d}\Qbb^M}{\mathrm{d}\Qbb^N}\Big|_{\F_t}
=
\frac{D(t)M(t)}{D(0)M(0)}
\frac{D(0)N(0)}{D(t)N(t)}
=
\frac{N(0)M(t)}{N(t)M(0)}.
\]

Hence, for every \(t\in[0,T]\),
\[
\boxed{
\frac{\mathrm{d}\Qbb^M}{\mathrm{d}\Qbb^N}\Big|_{\F_t}
=
\frac{N(0)M(t)}{N(t)M(0)}.
}
\]

In particular, at terminal time \(T\),
\[
\boxed{
\frac{\mathrm{d}\Qbb^M}{\mathrm{d}\Qbb^N}
=
\frac{N(0)M(T)}{N(T)M(0)}.
}
\]

\paragraph{Interpretation.}

This formula shows that the change of measure between two numeraires is
entirely determined by their relative value.

More precisely, changing numeraire from \(N\) to \(M\) amounts to
reweighting probabilities according to the ratio
\[
\frac{M(T)}{N(T)}.
\]

Thus the relation between the two associated measures is encoded in the
relative performance of the two numeraires.

\paragraph{A useful consequence.}

For every integrable random variable \(X\), Bayes' rule gives
\[
\E^{\Qbb^M}[X\mid\F_t]
=
\frac{1}{L_t^{M/N}}
\E^{\Qbb^N}[L_T^{M/N}X\mid\F_t],
\]
where
\[
L_t^{M/N}
:=
\frac{\mathrm{d}\Qbb^M}{\mathrm{d}\Qbb^N}\Big|_{\F_t}
=
\frac{N(0)M(t)}{N(t)M(0)}.
\]

This is the practical formula used to pass from one numeraire-based
pricing measure to another.

\paragraph{Interpretation.}
The density process \(L_t\) reweights probabilities according to the
relative performance of the chosen numeraire \(N\) compared with the
money market account.

If the numeraire performs well in some states of the world, those states
receive relatively more weight under \(\Qbb^N\).
Thus, changing numeraire amounts to changing both:
\begin{itemize}[leftmargin=1.4em]
\item the unit of account,
\item the probability measure.
\end{itemize}

\paragraph{Why equivalence matters.}
The measures \(\Qbb\) and \(\Qbb^N\) are equivalent, which means that
they assign zero probability to the same events.
Thus, changing numeraire does not change which events are possible; it
only changes their probabilities.

\paragraph{Main consequence.}
Under the measure \(\Qbb^N\), the price of any traded asset \(S(t)\)
expressed in units of the numeraire,
\[
\frac{S(t)}{N(t)},
\]
is a martingale.

This is the fundamental property that makes the change of numeraire so
useful in pricing.

\paragraph{Special case: money market account.}
If one chooses \(N(t)=B(t)\), then
\[
L_t = \frac{B(t)/B(t)}{B(0)/B(0)} = 1,
\]
so \(\Qbb^N=\Qbb\).
Hence the usual risk-neutral measure is simply the measure associated
with the money market account.

\paragraph{What we do next.}
In the next subsection, we state and prove the martingale property of
asset prices expressed in units of the chosen numeraire.


\subsection{Martingale property under the chosen numeraire}

We now state the key result behind the change of numeraire method.

\paragraph{Main statement.}
Let \(N(t)\) be a strictly positive traded asset chosen as numeraire, and
let \(\Qbb^N\) be the probability measure associated with \(N\).
Then, for any traded asset \(S(t)\), the denominated price process
\[
\frac{S(t)}{N(t)}
\]
is a \(\Qbb^N\)-martingale.

In other words, under the measure corresponding to the chosen
numeraire, all asset prices measured in units of that numeraire have
zero drift.

\paragraph{Interpretation.}
This is the exact analogue of the usual risk-neutral property.
Under the standard risk-neutral measure \(\Qbb\), discounted prices
\[
\frac{S(t)}{B(t)}
\]
are martingales, where \(B(t)\) is the money market account.
After changing numeraire from \(B\) to \(N\), the same property holds
with \(N\) in place of \(B\).

Thus the change of numeraire does not alter the fundamental no-arbitrage
principle; it only changes the unit in which prices are measured.

\paragraph{Conditional expectation form.}
Equivalently, for every \(0\le t\le T\),
\[
\frac{S(t)}{N(t)}
=
\E^{\Qbb^N}\!\left[
\frac{S(T)}{N(T)}
\;\middle|\;
\mathcal{F}_t
\right].
\]

This is the martingale representation that will lead directly to the
pricing formula under a general numeraire.

\paragraph{Why this is useful.}
The practical advantage is that one may choose a numeraire adapted to
the payoff being priced.
If the payoff naturally involves some asset \(N(T)\), then expressing
the claim in units of \(N(T)\) often simplifies the pricing formula and
the corresponding dynamics.

This is particularly useful:
\begin{itemize}[leftmargin=1.4em]
\item in fixed-income models, where bond numeraires lead to forward measures,
\item in FX models, where foreign and domestic money market accounts
naturally generate different pricing measures,
\item and more generally whenever the standard money-market discounting
is not the most convenient choice.
\end{itemize}

\paragraph{Heuristic proof.}
Recall that the measure \(\Qbb^N\) is defined by the density process
\[
L_t
=
\frac{D(t)N(t)}{D(0)N(0)},
\qquad D(t)=\frac1{B(t)}.
\]
Since both \(D(t)S(t)\) and \(D(t)N(t)\) are \(\Qbb\)-martingales, Bayes'
rule gives, for \(t\le T\),
\[
\E^{\Qbb^N}\!\left[\frac{S(T)}{N(T)}\middle|\mathcal{F}_t\right]
=
\frac{1}{L_t}
\E^\Qbb\!\left[L_T\frac{S(T)}{N(T)}\middle|\mathcal{F}_t\right].
\]
Using
\[
L_T\frac{S(T)}{N(T)}
=
\frac{D(T)N(T)}{D(0)N(0)}\frac{S(T)}{N(T)}
=
\frac{D(T)S(T)}{D(0)N(0)},
\]
we obtain
\[
\E^{\Qbb^N}\!\left[\frac{S(T)}{N(T)}\middle|\mathcal{F}_t\right]
=
\frac{1}{L_t}\,
\frac{1}{D(0)N(0)}
\E^\Qbb\!\left[D(T)S(T)\middle|\mathcal{F}_t\right].
\]
Since \(D(t)S(t)\) is a \(\Qbb\)-martingale,
\[
\E^\Qbb\!\left[D(T)S(T)\middle|\mathcal{F}_t\right]=D(t)S(t).
\]
Therefore
\[
\E^{\Qbb^N}\!\left[\frac{S(T)}{N(T)}\middle|\mathcal{F}_t\right]
=
\frac{D(t)S(t)}{D(t)N(t)}
=
\frac{S(t)}{N(t)},
\]
which proves the martingale property.

\paragraph{Key message.}
Once the numeraire \(N\) has been fixed, the correct martingale object
is no longer the discounted price \(D(t)S(t)\), but the denominated
price
\[
\frac{S(t)}{N(t)}.
\]
This is the fundamental principle behind all change-of-numeraire pricing
formulas.

\section{Pricing formula and change-of-measure intuition}

\subsection{Pricing under a general numeraire}

We now derive the pricing formula under an arbitrary numeraire.

\paragraph{Set-up.}
Let \(N(t)\) be a strictly positive traded asset chosen as numeraire, and
let \(\Qbb^N\) be the associated probability measure.
Consider a contingent claim paying \(X\) at maturity \(T\).

The time-\(t\) price of the claim will be denoted by \(\Pi(t;X)\).

\paragraph{Key idea.}
Since prices expressed in units of the numeraire are martingales under
\(\Qbb^N\), the denominated price process
\[
\frac{\Pi(t;X)}{N(t)}
\]
must satisfy
\[
\frac{\Pi(t;X)}{N(t)}
=
\E^{\Qbb^N}\!\left[
\frac{X}{N(T)}
\;\middle|\;
\mathcal{F}_t
\right].
\]

Multiplying both sides by \(N(t)\), we obtain the general pricing
formula
\[
\boxed{
\Pi(t;X)
=
N(t)\,
\E^{\Qbb^N}\!\left[
\frac{X}{N(T)}
\;\middle|\;
\mathcal{F}_t
\right].
}
\]

\paragraph{Interpretation.}
This is the analogue of the usual risk-neutral pricing formula, except
that the payoff is first expressed in units of the chosen numeraire,
and the expectation is taken under the corresponding measure.

Thus the pricing problem always has the same structure:
\begin{itemize}[leftmargin=1.4em]
\item divide the payoff by the numeraire at maturity,
\item take conditional expectation under the associated measure,
\item multiply back by the current value of the numeraire.
\end{itemize}

\paragraph{Comparison with the standard risk-neutral formula.}
If the numeraire is the money market account \(B(t)\), then
\[
\Pi(t;X)
=
B(t)\,
\E^\Qbb\!\left[
\frac{X}{B(T)}
\;\middle|\;
\mathcal{F}_t
\right].
\]
Since
\[
\frac{B(t)}{B(T)}=e^{-r(T-t)}
\]
in the constant-rate case, this reduces to the usual formula
\[
\Pi(t;X)
=
\E^\Qbb\!\left[
e^{-r(T-t)}X
\;\middle|\;
\mathcal{F}_t
\right].
\]

Hence the ordinary risk-neutral pricing formula is simply the special
case corresponding to the money market numeraire.

\paragraph{Why another numeraire may help.}
The usefulness of the formula
\[
\Pi(t;X)
=
N(t)\,
\E^{\Qbb^N}\!\left[
\frac{X}{N(T)}
\;\middle|\;
\mathcal{F}_t
\right]
\]
is that the choice of \(N\) may simplify the payoff or the dynamics.

For example:
\begin{itemize}[leftmargin=1.4em]
\item if the payoff naturally contains a bond \(P(T,T^\ast)\), choosing
that bond as numeraire may eliminate the discounting term;
\item if the payoff is naturally expressed relative to the stock price,
choosing the stock as numeraire may simplify the expectation;
\item in FX models, choosing the domestic or foreign money market account
leads to different but equivalent pricing representations.
\end{itemize}

\paragraph{A useful rule of thumb.}
The best numeraire is usually the one that makes the normalized payoff
\[
\frac{X}{N(T)}
\]
as simple as possible.

\paragraph{What comes next.}
In the following subsections, we illustrate this general pricing
principle with the most important examples:
the money market account, the stock as numeraire, and the zero-coupon
bond numeraire.

\subsection{Relation between numeraires and change of measure}

We now make explicit the relation between a change of numeraire and the
corresponding change of probability measure.

\paragraph{General principle.}
Choosing a new numeraire does not only modify the unit in which prices
are expressed; it also changes the probability measure under which the
denominated prices become martingales.

Thus, a change of numeraire always comes together with a change of
measure.

\paragraph{From the money market numeraire to a general numeraire.}
Let \(B(t)\) be the money market account, and let \(\Qbb\) be the usual
risk-neutral measure associated with \(B\).
Let \(N(t)\) be another strictly positive traded asset, used as a new
numeraire.

The measure \(\Qbb^N\) associated with \(N\) is defined through the
Radon--Nikodym density
\[
\frac{\mathrm{d}\Qbb^N}{\mathrm{d}\Qbb}\Big|_{\mathcal{F}_t}
=
\frac{N(t)/B(t)}{N(0)/B(0)}
=
\frac{D(t)N(t)}{D(0)N(0)},
\qquad D(t)=\frac1{B(t)}.
\]

Since \(D(t)N(t)\) is a positive \(\Qbb\)-martingale, this indeed defines
an equivalent probability measure.

\paragraph{Bayes' rule.}
For every integrable random variable \(X\), Bayes' rule gives
\[
\E^{\Qbb^N}[X\mid\mathcal{F}_t]
=
\frac{1}{L_t}\,
\E^\Qbb[L_T X\mid\mathcal{F}_t],
\qquad
L_t:=\frac{D(t)N(t)}{D(0)N(0)}.
\]

This formula is the practical tool that allows us to move from one
measure to another.

\paragraph{Interpretation.}
The density process \(L_t\) tilts probabilities according to the
relative performance of the new numeraire \(N\) with respect to the
money market account.

Hence, under the new measure \(\Qbb^N\), states of the world in which
the numeraire performs well receive relatively more weight.

\paragraph{Relation between two arbitrary numeraires.}
More generally, suppose that \(N^1(t)\) and \(N^2(t)\) are two strictly
positive traded numeraires, with associated measures \(\Qbb^{N^1}\) and
\(\Qbb^{N^2}\).
Then the change of measure from \(\Qbb^{N^1}\) to \(\Qbb^{N^2}\) is given by
\[
\frac{\mathrm{d}\Qbb^{N^2}}{\mathrm{d}\Qbb^{N^1}}\Big|_{\mathcal{F}_t}
=
\frac{N^2(t)/N^1(t)}{N^2(0)/N^1(0)}.
\]

Thus, the relative price process
\[
\frac{N^2(t)}{N^1(t)}
\]
plays the role of density process between the two associated measures.

\paragraph{Why this is natural.}
This formula shows that changing numeraire from \(N^1\) to \(N^2\) is
equivalent to reweighting probabilities by the relative value of the two
numeraires.

In particular, the change of numeraire is entirely encoded in the ratio
of the two numeraires.

\paragraph{Main consequence.}
Once the measure has been changed accordingly, any traded asset \(S(t)\)
expressed in units of the new numeraire becomes a martingale:
\[
\frac{S(t)}{N^2(t)}
\quad\text{is a }\Qbb^{N^2}\text{-martingale.}
\]

This is the key mechanism behind all practical applications of the
change-of-numeraire method.

\paragraph{What we do next.}
To make this more concrete, we next revisit It\^o calculus for ratio
processes and see how the dynamics of asset prices change when the
numeraire is replaced.



\section{Examples and applications}


\subsection{Money market account}

We start with the most classical choice of numeraire: the money market
account.

\paragraph{Definition.}
The money market account \(B(t)\) is the value at time \(t\) of one unit
invested in the bank account at time \(0\).
Its dynamics are
\[
\mathrm{d}B(t)=R(t)B(t)\,\mathrm{d}t,
\qquad B(0)=1,
\]
so that
\[
B(t)=\exp\!\Big(\int_0^t R(s)\,\mathrm{d}s\Big).
\]
In the constant-rate case, this reduces to
\[
B(t)=e^{rt}.
\]

\paragraph{Associated numeraire.}
If we choose \(B(t)\) as numeraire, then the denominated price of an
asset \(S(t)\) is simply its discounted price
\[
\frac{S(t)}{B(t)}=D(t)S(t),
\qquad D(t):=\frac1{B(t)}.
\]

\paragraph{Associated measure.}
The probability measure associated with the money market account is the
usual risk-neutral measure, which we denote by \(\Qbb\).
Under this measure, every traded asset expressed in units of the money
market account is a martingale.

In other words, for any traded asset \(S(t)\),
\[
\frac{S(t)}{B(t)}
\quad\text{is a }\Qbb\text{-martingale.}
\]

Equivalently, for \(0\le t\le T\),
\[
\frac{S(t)}{B(t)}
=
\E^\Qbb\!\left[
\frac{S(T)}{B(T)}
\;\middle|\;
\mathcal{F}_t
\right].
\]

\paragraph{Pricing formula.}
For a contingent claim paying \(X\) at time \(T\), the general
numeraire-pricing formula becomes
\[
\Pi(t;X)
=
B(t)\E^\Qbb\!\left[
\frac{X}{B(T)}
\;\middle|\;
\mathcal{F}_t
\right].
\]
In the constant-rate case,
\[
\boxed{
\Pi(t;X)
=
\E^\Qbb\!\left[
e^{-r(T-t)}X
\;\middle|\;
\mathcal{F}_t
\right].
}
\]

This is the standard risk-neutral pricing formula used throughout the
course.

\paragraph{Dynamics under the risk-neutral measure.}
If the asset price follows
\[
\mathrm{d}S(t)=\mu(t)S(t)\,\mathrm{d}t+\sigma(t)S(t)\,\mathrm{d}W(t)
\]
under some measure, then under the risk-neutral measure \(\Qbb\) its
drift becomes the short rate:
\[
\mathrm{d}S(t)=R(t)S(t)\,\mathrm{d}t+\sigma(t)S(t)\,\mathrm{d}W^\Qbb(t).
\]
Hence
\[
\mathrm{d}\Big(\frac{S(t)}{B(t)}\Big)
=
\frac{S(t)}{B(t)}\sigma(t)\,\mathrm{d}W^\Qbb(t),
\]
which makes the martingale property explicit.

\paragraph{Interpretation.}
Choosing the money market account as numeraire means that all prices are
measured in units of cash invested at the short rate.
This is the most natural choice when interest rates are deterministic or
when no other numeraire simplifies the problem.

\paragraph{Why this is the standard choice.}
The money market account is the default numeraire because:
\begin{itemize}[leftmargin=1.4em]
\item it leads to the usual risk-neutral valuation formula,
\item discounting by \(B(t)\) is familiar and easy to interpret,
\item it is the natural benchmark in equity option pricing.
\end{itemize}

\paragraph{Limitation.}
Although the money market numeraire is the standard one, it is not
always the most convenient.
For some payoffs, especially in fixed-income or FX models, another
choice of numeraire leads to a simpler dynamics or a cleaner pricing
formula.
This is precisely why the change-of-numeraire method is so useful.


\subsection{Stock as numeraire}

We now consider another important choice of numeraire: the stock itself.

\paragraph{Definition.}
Suppose that the stock price \(S(t)\) is strictly positive for all
\(t\in[0,T]\).
Then we may choose \(S(t)\) as numeraire.
In that case, the denominated price of any traded asset \(X(t)\) is
\[
\frac{X(t)}{S(t)}.
\]

\paragraph{Associated measure.}
The probability measure associated with the stock numeraire will be
denoted by \(\Qbb^S\).
By the general change-of-numeraire principle, every traded asset
expressed in units of the stock is a \(\Qbb^S\)-martingale.

Thus, for any traded asset \(X(t)\),
\[
\frac{X(t)}{S(t)}
\quad\text{is a }\Qbb^S\text{-martingale.}
\]

Equivalently, for \(0\le t\le T\),
\[
\frac{X(t)}{S(t)}
=
\E^{\Qbb^S}\!\left[
\frac{X(T)}{S(T)}
\;\middle|\;
\mathcal{F}_t
\right].
\]

\paragraph{Pricing formula.}
Let \(X\) be a contingent claim paid at time \(T\).
The general pricing formula becomes
\[
\Pi(t;X)
=
S(t)\E^{\Qbb^S}\!\left[
\frac{X}{S(T)}
\;\middle|\;
\mathcal{F}_t
\right].
\]

This representation is especially convenient when the payoff naturally
contains a factor \(S(T)\).

\paragraph{Typical example.}
Suppose that the claim payoff can be written as
\[
X=S(T)\,Y,
\]
for some random variable \(Y\).
Then the stock-numeraire pricing formula simplifies to
\[
\Pi(t;X)
=
S(t)\E^{\Qbb^S}\!\left[
Y
\;\middle|\;
\mathcal{F}_t
\right].
\]

So the factor \(S(T)\) disappears from the expectation once the stock is
used as numeraire.

\paragraph{Connection with equity derivatives.}
This choice of numeraire is useful for payoffs that are naturally
homogeneous in the stock price.
It also provides an elegant alternative representation for some
equity-linked contracts.

In particular, if a payoff is proportional to the terminal stock price,
then working under \(\Qbb^S\) often leads to a simpler expectation than
under the usual risk-neutral measure.

\paragraph{Density process.}
Starting from the usual risk-neutral measure \(\Qbb\), the stock
numeraire measure \(\Qbb^S\) is defined by
\[
\frac{\mathrm{d}\Qbb^S}{\mathrm{d}\Qbb}\Big|_{\mathcal{F}_t}
=
\frac{D(t)S(t)}{D(0)S(0)},
\qquad D(t)=\frac1{B(t)}.
\]

Since \(D(t)S(t)\) is a positive \(\Qbb\)-martingale, this defines an
equivalent probability measure.

\paragraph{Interpretation.}
Under \(\Qbb^S\), probabilities are tilted according to the performance
of the stock relative to the money market account.
States of the world in which the stock performs well receive relatively
more weight.

\paragraph{Why this matters.}
The stock numeraire is less common than the money market account in
basic pricing formulas, but it is conceptually important:
it shows that the same derivative can be priced under different
measures, depending on the chosen numeraire.

This flexibility is one of the main strengths of the
change-of-numeraire method.


\subsection{Zero-coupon bond as numeraire}

A particularly important choice of numeraire in fixed-income modeling is
a zero-coupon bond.

\paragraph{Definition.}
Fix a maturity date \(T^\ast\).
We denote by
\[
P(t,T^\ast)
\]
the price at time \(t\leq T^\ast\) of a zero-coupon bond paying \(1\) at
time \(T^\ast\).

Since bond prices are strictly positive, \(P(t,T^\ast)\) can be used as
a numeraire.

\paragraph{Associated measure.}
The probability measure associated with the numeraire
\[
N(t)=P(t,T^\ast)
\]
is called the \emph{\(T^\ast\)-forward measure}, and is usually denoted
by
\[
\Qbb^{T^\ast}.
\]

By the general change-of-numeraire principle, any traded asset price
expressed in units of \(P(t,T^\ast)\) is a \(\Qbb^{T^\ast}\)-martingale.

Thus, for any traded asset \(X(t)\),
\[
\frac{X(t)}{P(t,T^\ast)}
\quad\text{is a }\Qbb^{T^\ast}\text{-martingale.}
\]

Equivalently,
\[
\frac{X(t)}{P(t,T^\ast)}
=
\E^{\Qbb^{T^\ast}}\!\left[
\frac{X(T)}{P(T,T^\ast)}
\;\middle|\;
\mathcal{F}_t
\right].
\]

\paragraph{Pricing formula.}
Let \(X\) be a contingent claim paid at time \(T\leq T^\ast\).
Then
\[
\Pi(t;X)
=
P(t,T^\ast)\,
\E^{\Qbb^{T^\ast}}\!\left[
\frac{X}{P(T,T^\ast)}
\;\middle|\;
\mathcal{F}_t
\right].
\]

In the important special case \(T=T^\ast\), we have
\[
P(T^\ast,T^\ast)=1,
\]
so the pricing formula simplifies to
\[
\boxed{
\Pi(t;X)
=
P(t,T^\ast)\,
\E^{\Qbb^{T^\ast}}\!\left[
X
\;\middle|\;
\mathcal{F}_t
\right].
}
\]


\paragraph{A useful Black--Scholes interpretation under the \(T\)-forward measure.}

One particularly useful consequence of the bond-numeraire approach is
the following simple rule of thumb:

\begin{quote}
under the \(T\)-forward measure, one may price a claim maturing at time
\(T\) as if the short rate were equal to \(0\), and then multiply the
result by \(P(0,T)\).
\end{quote}

\paragraph{Why this is true.}

Suppose that a contingent claim pays \(X\) at time \(T\).
Using the zero-coupon bond \(P(t,T)\) as numeraire, the pricing formula
becomes
\[
\Pi(t;X)
=
P(t,T)\,
\E^{\Qbb^T}\!\left[
X
\;\middle|\;
\F_t
\right],
\]
since
\[
P(T,T)=1.
\]

Thus the discounting disappears from the expectation.
In this sense, the role played by the short rate in the usual
risk-neutral formula is absorbed into the bond numeraire.

\paragraph{Forward price viewpoint.}

If \(S(t)\) is a traded asset, define its \(T\)-forward price by
\[
F(t,T):=\frac{S(t)}{P(t,T)}.
\]

Under the \(T\)-forward measure, this normalized price process is a
martingale.
Equivalently, one may think of \(F(t,T)\) as evolving like an asset with
zero drift under the corresponding pricing measure.

This is exactly the analogue of the Black--Scholes setting with
vanishing interest rate.

\paragraph{Practical rule.}

Suppose that the payoff at time \(T\) is of the form
\[
X=h(S(T)).
\]
Then, instead of working directly under the money-market numeraire, one
may:
\begin{itemize}[leftmargin=1.4em]
\item switch to the \(T\)-forward measure,
\item work with the forward price process,
\item use the Black--Scholes formula with effective rate \(r=0\),
\item and finally multiply the result by \(P(0,T)\).
\end{itemize}

In other words, if \(C^{(0)}(0)\) denotes the Black--Scholes price
computed with zero short rate and forward-style normalization, then
\[
\boxed{
\Pi(0;X)=P(0,T)\,C^{(0)}(0).
}
\]

\paragraph{Interpretation.}

The point is not that the interest rate has disappeared from the model,
but rather that it has been absorbed into the numeraire.
The pricing problem is therefore rewritten in a form where the relevant
normalized process has zero drift under the new measure.

\paragraph{Pedagogical message.}

The \(T\)-forward measure is often useful because it transforms pricing
problems with discounting into pricing problems without explicit
discounting inside the expectation.

This is why one often says informally that, under the \(T\)-forward
measure, ``Black--Scholes works with \(r=0\)'', provided one remembers to
multiply the final answer by the bond price \(P(0,T)\).


\paragraph{Why this is useful.}
This is the main reason why bond numeraires are so important.
If the payoff is received exactly at the maturity of the chosen bond,
the discounting disappears from the expectation.
The price is simply the bond price times a conditional expectation under
the corresponding forward measure.

\paragraph{Forward prices.}
A natural quantity in fixed-income markets is the forward price of an
asset for delivery at time \(T^\ast\).
Under the \(T^\ast\)-forward measure, such forward prices typically
become martingales.
This is completely analogous to the usual risk-neutral framework, except
that the bond \(P(t,T^\ast)\) replaces the money market account.

\paragraph{Density process.}
Starting from the usual risk-neutral measure \(\Qbb\), the
\(T^\ast\)-forward measure is defined by
\[
\frac{\mathrm{d}\Qbb^{T^\ast}}{\mathrm{d}\Qbb}\Big|_{\mathcal{F}_t}
=
\frac{D(t)P(t,T^\ast)}{D(0)P(0,T^\ast)},
\qquad D(t)=\frac1{B(t)}.
\]

Since \(D(t)P(t,T^\ast)\) is a positive \(\Qbb\)-martingale, this indeed
defines an equivalent probability measure.

\paragraph{Interpretation.}
Using a zero-coupon bond as numeraire means that prices are measured in
units of future cash delivered at time \(T^\ast\).
This is particularly natural in interest-rate models, where many claims
are linked to future bond prices, forward rates, or swap rates.

\paragraph{Main application.}
The bond numeraire and the associated forward measure are fundamental
tools in fixed-income derivative pricing.
They simplify the pricing of bond options, caplets, floorlets, and many
other interest-rate products.

\paragraph{Key message.}
Choosing a zero-coupon bond as numeraire is often the most natural
choice when the payoff is linked to a future date \(T^\ast\).
In that case, the corresponding forward measure leads to simpler pricing
formulas and more tractable dynamics.


\subsection{Forward measures}

Forward measures arise naturally when the chosen numeraire is a
zero-coupon bond.

\paragraph{Definition.}
Fix a maturity date \(T^\ast\), and consider the zero-coupon bond
\[
P(t,T^\ast).
\]
If this bond is used as numeraire, the associated probability measure is
called the \emph{\(T^\ast\)-forward measure}, and is denoted by
\[
\Qbb^{T^\ast}.
\]

Thus the \(T^\ast\)-forward measure is simply the measure associated
with the bond numeraire \(P(t,T^\ast)\).

\paragraph{Key martingale property.}
By the general change-of-numeraire principle, any traded asset price
expressed in units of \(P(t,T^\ast)\) is a \(\Qbb^{T^\ast}\)-martingale.

In particular, if \(X\) is a contingent claim paid at time \(T^\ast\),
then its price process satisfies
\[
\frac{\Pi(t;X)}{P(t,T^\ast)}
=
\E^{\Qbb^{T^\ast}}\!\left[
X
\;\middle|\;
\mathcal{F}_t
\right].
\]

Equivalently,
\[
\boxed{
\Pi(t;X)
=
P(t,T^\ast)\,
\E^{\Qbb^{T^\ast}}\!\left[
X
\;\middle|\;
\mathcal{F}_t
\right].
}
\]

This is the basic pricing formula under the forward measure.

\paragraph{Why this is useful.}
The forward measure is particularly convenient when the payoff is paid at
a fixed future date \(T^\ast\).
In that case, the discounting disappears from the expectation, and the
pricing formula becomes especially simple.

This is one of the main reasons why forward measures play such an
important role in fixed-income modeling.

\paragraph{Forward prices.}
Suppose that an asset \(X(t)\) can be delivered at time \(T^\ast\).
Its forward price for delivery at time \(T^\ast\) is typically defined by
normalizing with the bond price \(P(t,T^\ast)\).

Under the \(T^\ast\)-forward measure, this forward price process is a
martingale.
This is the natural analogue of the usual risk-neutral martingale
property, but with the bond \(P(t,T^\ast)\) replacing the money market
account.

\paragraph{Example: bond option intuition.}
Suppose that a claim pays
\[
X = h\big(P(T^\ast,T_1),\dots,P(T^\ast,T_n)\big)
\]
at time \(T^\ast\).
Then under the \(T^\ast\)-forward measure, its price is
\[
\Pi(t;X)
=
P(t,T^\ast)\,
\E^{\Qbb^{T^\ast}}\!\left[
h\big(P(T^\ast,T_1),\dots,P(T^\ast,T_n)\big)
\;\middle|\;
\mathcal{F}_t
\right].
\]

Thus the bond price \(P(t,T^\ast)\) factors out, and the remaining
expectation is taken under a measure naturally adapted to the payment
date.

\paragraph{Connection with LIBOR and swap market models.}
In practical interest-rate modeling, different forward measures are
attached to different payment dates.
For example:
\begin{itemize}[leftmargin=1.4em]
\item caplets are naturally priced under the forward measure associated
with their payment date,
\item swap-related quantities are often studied under annuity measures,
which are built from suitable linear combinations of bonds.
\end{itemize}

Thus forward measures are one of the main building blocks of modern
fixed-income derivative pricing.

\paragraph{Density process.}
Starting from the standard risk-neutral measure \(\Qbb\), the
\(T^\ast\)-forward measure is defined by
\[
\frac{\mathrm{d}\Qbb^{T^\ast}}{\mathrm{d}\Qbb}\Big|_{\mathcal{F}_t}
=
\frac{D(t)P(t,T^\ast)}{D(0)P(0,T^\ast)},
\qquad D(t)=\frac1{B(t)}.
\]

Since \(D(t)P(t,T^\ast)\) is a positive \(\Qbb\)-martingale, this indeed
defines an equivalent probability measure.

\paragraph{Interpretation.}
The forward measure is the probability measure under which future cash
flows paid at time \(T^\ast\) are naturally valued in units of one unit
of currency delivered at time \(T^\ast\).

\paragraph{Key message.}
Forward measures are not new pricing principles; they are simply the
change-of-numeraire principle applied to zero-coupon bonds.
Their importance comes from the fact that they are often the most
natural and efficient measures for pricing fixed-income derivatives.


\subsection{Applications to derivative pricing}

We now illustrate how the change-of-numeraire method can be used in
practice to simplify derivative pricing formulas.

\paragraph{General strategy.}
Suppose that a contingent claim pays \(X\) at time \(T\).
Under a numeraire \(N\) with associated measure \(\Qbb^N\), its price is
given by
\[
\Pi(t;X)
=
N(t)\,
\E^{\Qbb^N}\!\left[
\frac{X}{N(T)}
\;\middle|\;
\mathcal{F}_t
\right].
\]

The practical idea is therefore simple:
\begin{itemize}[leftmargin=1.4em]
\item choose a numeraire adapted to the structure of the payoff,
\item rewrite the payoff in units of this numeraire,
\item compute the expectation under the associated measure.
\end{itemize}

A good choice of numeraire often leads to a much simpler pricing
formula.

\paragraph{Application 1: standard risk-neutral pricing.}
If the money market account \(B(t)\) is chosen as numeraire, then we
recover the usual risk-neutral formula
\[
\Pi(t;X)
=
B(t)\E^\Qbb\!\left[
\frac{X}{B(T)}
\;\middle|\;
\mathcal{F}_t
\right].
\]
In the constant-rate case, this becomes
\[
\Pi(t;X)
=
\E^\Qbb\!\left[
e^{-r(T-t)}X
\;\middle|\;
\mathcal{F}_t
\right].
\]

This is the default pricing formula used in most equity-option models.

\paragraph{Application 2: stock numeraire.}
Suppose that the payoff can be written as
\[
X=S(T)\,Y,
\]
for some random variable \(Y\).
Then using the stock as numeraire gives
\[
\Pi(t;X)
=
S(t)\E^{\Qbb^S}\!\left[
Y
\;\middle|\;
\mathcal{F}_t
\right].
\]

Thus the factor \(S(T)\) disappears inside the expectation.
This can be useful for equity-linked contracts whose payoff is naturally
proportional to the stock price.

\paragraph{Application 3: bond numeraire.}
Suppose now that the payoff is paid at a fixed date \(T^\ast\).
Choosing the zero-coupon bond \(P(t,T^\ast)\) as numeraire yields
\[
\Pi(t;X)
=
P(t,T^\ast)\,
\E^{\Qbb^{T^\ast}}\!\left[
X
\;\middle|\;
\mathcal{F}_t
\right].
\]

This is particularly convenient for fixed-income derivatives, because
the discounting disappears from the expectation.

For example, if a claim pays a function of bond prices at time
\(T^\ast\), say
\[
X=h\big(P(T^\ast,T_1),\dots,P(T^\ast,T_n)\big),
\]
then its price becomes
\[
\Pi(t;X)
=
P(t,T^\ast)\,
\E^{\Qbb^{T^\ast}}\!\left[
h\big(P(T^\ast,T_1),\dots,P(T^\ast,T_n)\big)
\;\middle|\;
\mathcal{F}_t
\right].
\]

\paragraph{Application 4: forward contracts.}
The change-of-numeraire method also explains naturally why forward
prices are martingales under the appropriate forward measure.

If delivery takes place at time \(T^\ast\), then under the
\(T^\ast\)-forward measure, forward prices are naturally expressed in
units of \(P(t,T^\ast)\), and therefore become martingales.

This is one of the central ideas in fixed-income modeling.

\paragraph{Application 5: foreign exchange intuition.}
In FX markets, one may choose either the domestic or the foreign money
market account as numeraire.
This leads to different but equivalent pricing measures.
The resulting formulas provide a natural framework for pricing foreign
exchange derivatives and for understanding the symmetry between
domestic and foreign viewpoints.

\paragraph{Main lesson.}
The change-of-numeraire method is not a new pricing principle; it is a
way of reorganizing the same no-arbitrage pricing problem so that the
payoff and the dynamics are expressed under the most convenient unit of
account.

\paragraph{Rule of thumb.}
Whenever a derivative pricing formula looks complicated, a natural
question to ask is:
\begin{quote}
Is there a better numeraire under which the payoff becomes simpler?
\end{quote}

This is often the key step that transforms a difficult pricing problem
into a tractable one.


% ============================================================
\section{Foreign exchange derivative pricing}
% ============================================================

We conclude with a brief application of the change-of-numeraire method
to foreign exchange markets.

The main idea is that an exchange rate can naturally be viewed from two
different perspectives:
\begin{itemize}[leftmargin=1.4em]
\item as the price of one unit of foreign currency expressed in domestic currency,
\item or, equivalently, through its inverse, as the price of one unit of domestic currency expressed in foreign currency.
\end{itemize}

This symmetry makes FX markets a particularly natural setting for the
change-of-numeraire method.

% ------------------------------------------------------------
\subsection{Domestic and foreign money market accounts}

We denote by
\[
B_d(t)=e^{r_d t}
\qquad\text{and}\qquad
B_f(t)=e^{r_f t}
\]
the domestic and foreign money market accounts, where \(r_d\) and \(r_f\)
are the domestic and foreign short rates, assumed here constant for
simplicity.

We denote by
\[
X(t)
\]
the spot exchange rate, defined as the price in domestic currency of one
unit of foreign currency.

Thus:
\begin{itemize}[leftmargin=1.4em]
\item \(X(t)\) is measured in domestic currency,
\item one unit of foreign cash invested in the foreign bank account has
domestic value
\[
X(t)B_f(t).
\]
\end{itemize}

This quantity plays a central role in FX pricing.

% ------------------------------------------------------------
\subsection{Risk-neutral dynamics of the exchange rate}

Under the domestic risk-neutral measure \(\Qbb^d\), the domestic
discounted value of any traded asset must be a martingale.

In particular, the foreign bank account expressed in domestic currency,
namely
\[
X(t)B_f(t),
\]
must satisfy
\[
\frac{X(t)B_f(t)}{B_d(t)}
\quad\text{is a }\Qbb^d\text{-martingale.}
\]

This implies that the exchange rate dynamics under the domestic
risk-neutral measure take the form
\[
\boxed{
\mathrm{d}X(t)=\big(r_d-r_f\big)X(t)\,\mathrm{d}t+\sigma X(t)\,\mathrm{d}W_t^{\Qbb^d}.
}
\]

\paragraph{Interpretation.}
The foreign interest rate plays the same role as a continuous dividend
yield.
Indeed, holding one unit of foreign currency gives access to the foreign
money market account, just as holding a stock with dividends provides a
continuous cashflow.

For this reason, an FX rate behaves mathematically like an equity with
continuous dividend yield \(r_f\).

% ------------------------------------------------------------
\subsection{Pricing formula under the domestic measure}

Let \(H(X(T))\) be the payoff of a European FX derivative, paid at time
\(T\) in domestic currency.

Under the domestic money market numeraire, its price is
\[
\boxed{
\Pi(t)
=
\E^{\Qbb^d}\!\left[
e^{-r_d(T-t)}H(X(T))
\;\middle|\;
\mathcal{F}_t
\right].
}
\]

This is the usual domestic risk-neutral pricing formula.

In particular, a European call option on the exchange rate with strike
\(K\) has payoff
\[
(X(T)-K)^+,
\]
and price
\[
\Pi(t)
=
\E^{\Qbb^d}\!\left[
e^{-r_d(T-t)}(X(T)-K)^+
\;\middle|\;
\mathcal{F}_t
\right].
\]

% ------------------------------------------------------------
\subsection{Black--Scholes formula for FX options}

Since the FX rate follows a lognormal dynamics under \(\Qbb^d\), the
usual Black--Scholes formula applies, with the foreign rate \(r_f\)
playing the role of a dividend yield.

Hence the domestic price of a European call option on the exchange rate
is
\[
\boxed{
C(t)
=
X(t)e^{-r_f(T-t)}N(d_1)
-
Ke^{-r_d(T-t)}N(d_2),
}
\]
where
\[
d_1
=
\frac{\ln(X(t)/K)+\big(r_d-r_f+\tfrac12\sigma^2\big)(T-t)}
{\sigma\sqrt{T-t}},
\]
\[
d_2=d_1-\sigma\sqrt{T-t}.
\]

Similarly, the European put price is
\[
\boxed{
P(t)
=
Ke^{-r_d(T-t)}N(-d_2)
-
X(t)e^{-r_f(T-t)}N(-d_1).
}
\]

\paragraph{Interpretation.}
This is exactly the Black--Scholes formula with:
\begin{itemize}[leftmargin=1.4em]
\item spot price \(X(t)\),
\item domestic discounting at rate \(r_d\),
\item foreign rate \(r_f\) acting as a continuous yield.
\end{itemize}

% ------------------------------------------------------------
\subsection{Change of numeraire viewpoint}

The FX setting provides a natural illustration of the change of
numeraire idea.

From the domestic viewpoint, we discount by \(B_d(t)\) and work under
the domestic risk-neutral measure \(\Qbb^d\).

From the foreign viewpoint, one may instead use the foreign money market
account \(B_f(t)\) as numeraire, or equivalently work with domestic
asset prices expressed in foreign currency.
This leads to a foreign pricing measure \(\Qbb^f\).

Thus FX pricing naturally involves two equivalent martingale measures:
\begin{itemize}[leftmargin=1.4em]
\item the domestic measure,
\item the foreign measure.
\end{itemize}

The same contract may then be represented in two different but
equivalent ways, depending on the chosen numeraire.

\paragraph{Main lesson.}
FX pricing is one of the clearest examples showing that:
\begin{quote}
the choice of numeraire is not just a mathematical trick, but a natural
way of reflecting different economic viewpoints.
\end{quote}

% ------------------------------------------------------------
\subsection{FX put--call parity}

Let \(C(t)\) and \(P(t)\) be the domestic prices of a European call and
put on the exchange rate with strike \(K\) and maturity \(T\).

The put--call parity relation becomes
\[
\boxed{
C(t)-P(t)=X(t)e^{-r_f(T-t)}-Ke^{-r_d(T-t)}.
}
\]

This is the FX analogue of the usual equity parity relation, with the
foreign rate \(r_f\) replacing the dividend yield.

% ------------------------------------------------------------
\subsection{Summary}

The key formulas to remember are:
\begin{itemize}[leftmargin=1.4em]
\item under the domestic measure,
\[
\mathrm{d}X(t)=\big(r_d-r_f\big)X(t)\,\mathrm{d}t+\sigma X(t)\,\mathrm{d}W_t^{\Qbb^d},
\]
\item the foreign rate acts like a dividend yield,
\item FX options are priced by Black--Scholes with \(r_d\) as discount
rate and \(r_f\) as yield,
\item the domestic and foreign viewpoints correspond to different
numeraires and different but equivalent martingale measures.
\end{itemize}

\paragraph{Domestic and foreign measures in FX markets.}

Let \(B_d(t)\) and \(B_f(t)\) denote the domestic and foreign money
market accounts, and let \(X(t)\) be the exchange rate, defined as the
price in domestic currency of one unit of foreign currency.

Under the domestic viewpoint, the natural numeraire is \(B_d(t)\), and
the associated pricing measure is the domestic risk-neutral measure
\(\Qbb^d\).
Under the foreign viewpoint, the natural numeraire is \(B_f(t)\), and
the associated pricing measure is the foreign risk-neutral measure
\(\Qbb^f\).

However, since \(X(t)\) converts foreign currency into domestic
currency, the foreign bank account expressed in domestic currency is
\[
X(t)B_f(t).
\]
This is a strictly positive traded asset, and it can therefore be used
as a domestic numeraire.

\paragraph{Change of measure between domestic and foreign viewpoints.}

Using the general numeraire-change formula, the Radon--Nikodym density
between the foreign and domestic measures is given by
\[
\frac{\mathrm{d}\Qbb^f}{\mathrm{d}\Qbb^d}\Big|_{\F_t}
=
\frac{B_d(0)\,X(t)B_f(t)}{B_d(t)\,X(0)B_f(0)}.
\]

In the constant-rate case, where
\[
B_d(t)=e^{r_d t},
\qquad
B_f(t)=e^{r_f t},
\]
this becomes
\[
\frac{\mathrm{d}\Qbb^f}{\mathrm{d}\Qbb^d}\Big|_{\F_t}
=
\frac{X(t)e^{r_f t}}{X(0)e^{r_d t}}.
\]

Equivalently, at time \(T\),
\[
\boxed{
\frac{\mathrm{d}\Qbb^f}{\mathrm{d}\Qbb^d}
=
\frac{X(T)e^{r_f T}}{X(0)e^{r_d T}}.
}
\]

\paragraph{Interpretation.}

This formula shows that moving from the domestic measure to the foreign
measure amounts to reweighting probabilities according to the relative
performance of the foreign bank account expressed in domestic currency.

Thus the domestic and foreign pricing measures are not unrelated:
they are linked exactly by the exchange rate and the two money market
accounts.

\paragraph{FX forward pricing.}

Consider now a forward contract written on the exchange rate, with
maturity \(T\) and delivery price \(K\).
Its payoff at time \(T\), in domestic currency, is
$$
X(T)-K.
$$

Its time-\(t\) price under the domestic measure is therefore
$$
\Pi(t)
=
\E^{\Qbb^d}\!\left[
e^{-r_d(T-t)}(X(T)-K)
\;\middle|\;
\F_t
\right].
$$

Using the fact that under \(\Qbb^d\),
$$
\mathrm{d}X(t)=\big(r_d-r_f\big)X(t)\,\mathrm{d}t+\sigma X(t)\,\mathrm{d}W_t^{\Qbb^d},
$$
we obtain
$$
\E^{\Qbb^d}[X(T)\mid\F_t]
=
X(t)e^{(r_d-r_f)(T-t)}.
$$

Hence
$$
\Pi(t)
=
e^{-r_d(T-t)}X(t)e^{(r_d-r_f)(T-t)}
-
e^{-r_d(T-t)}K,
$$
that is,
$$
\boxed{
\Pi(t)=X(t)e^{-r_f(T-t)}-Ke^{-r_d(T-t)}.
}
$$

\paragraph{Forward exchange rate.}

The forward exchange rate \(F(t,T)\) is defined by the value of \(K\)
for which the contract has zero value at time \(t\).
Thus
$$
\Pi(t)=0
\qquad\Longleftrightarrow\qquad
K=F(t,T),
$$
so that
$$
\boxed{
F(t,T)=X(t)e^{(r_d-r_f)(T-t)}.
}
$$

\paragraph{Interpretation.}

This is the FX analogue of the usual forward-pricing formula.
The domestic rate pushes the forward price upward, while the foreign
rate acts like a continuous dividend yield and pushes it downward.

\paragraph{Key message.}

The FX market provides a particularly transparent example of the
change-of-numeraire method:
\begin{itemize}[leftmargin=1.4em]
\item the domestic and foreign viewpoints correspond to different
numeraires,
\item the corresponding pricing measures are linked by an explicit
Radon--Nikodym density,
\item and the forward pricing formula follows naturally from the
domestic risk-neutral dynamics.
\end{itemize}


% ============================================================
\section{Interpretation and summary}
% ============================================================

\subsection{Link with risk-neutral valuation}

The change-of-numeraire method should be understood as a generalization
of the usual risk-neutral valuation principle.

In the standard framework, the money market account \(B(t)\) is chosen
as numeraire, and the associated probability measure is the usual
risk-neutral measure \(\Qbb\).
Under this measure, discounted asset prices
\[
\frac{S(t)}{B(t)}
\]
are martingales, and the price of a contingent claim \(X\) is given by
\[
\Pi(t;X)
=
B(t)\E^\Qbb\!\left[
\frac{X}{B(T)}
\;\middle|\;
\mathcal{F}_t
\right].
\]

The change-of-numeraire method extends this idea by replacing the money
market account with an arbitrary strictly positive traded asset \(N(t)\).

Once \(N\) is chosen as numeraire, the relevant martingale object is no
longer the discounted price \(S(t)/B(t)\), but the denominated price
\[
\frac{S(t)}{N(t)}.
\]
The corresponding pricing measure is no longer \(\Qbb\), but the
measure \(\Qbb^N\) associated with \(N\).

The pricing formula then becomes
\[
\Pi(t;X)
=
N(t)\E^{\Qbb^N}\!\left[
\frac{X}{N(T)}
\;\middle|\;
\mathcal{F}_t
\right].
\]

Thus the usual risk-neutral valuation formula is simply the special case
corresponding to the money market numeraire.

\paragraph{Key message.}
Risk-neutral valuation is not tied to one specific measure.
Rather, it is a general no-arbitrage pricing principle:
\begin{quote}
once a numeraire is chosen, there exists an associated equivalent
martingale measure under which prices expressed in units of that
numeraire are martingales.
\end{quote}

In this sense, the ordinary risk-neutral measure is only one member of a
whole family of equivalent martingale measures indexed by the choice of
numeraire.

\subsection{Practical interpretation}

From a practical point of view, the main benefit of the
change-of-numeraire method is that it provides flexibility.

The no-arbitrage price of a contingent claim is unique, but the way we
represent that price may vary depending on the numeraire and the
associated measure.
A good choice of numeraire can make both the pricing formula and the
underlying dynamics much simpler.

\paragraph{How to think about it in practice.}
When pricing a derivative, one may proceed as follows:
\begin{itemize}[leftmargin=1.4em]
\item identify the payoff date and the structure of the payoff,
\item look for a strictly positive traded asset naturally related to
that payoff,
\item use this asset as numeraire,
\item express the payoff in units of the chosen numeraire,
\item compute the expectation under the associated measure.
\end{itemize}

This often leads to a much cleaner formula than under the standard
money-market numeraire.

\paragraph{Typical situations.}
\begin{itemize}[leftmargin=1.4em]
\item In equity pricing, the money market account is often the most
natural choice.
\item In fixed-income pricing, a zero-coupon bond maturing at the payoff
date typically leads to a forward measure and a simpler formula.
\item In FX modeling, domestic and foreign money-market accounts lead to
different but equivalent pricing measures.
\end{itemize}

\paragraph{Conceptual summary.}
The change-of-numeraire method does not change the economic content of
pricing.
It simply changes:
\begin{itemize}[leftmargin=1.4em]
\item the unit in which prices are measured,
\item the probability measure under which the martingale property holds.
\end{itemize}

The price itself remains the same; only its representation changes.

\paragraph{Final takeaway.}
The change of numeraire is one of the most useful tools in modern
derivative pricing.
It provides:
\begin{itemize}[leftmargin=1.4em]
\item a deeper understanding of risk-neutral valuation,
\item a unifying framework for different pricing measures,
\item and a practical method for simplifying derivative pricing
problems.
\end{itemize}


\end{document}